\documentclass[conference]{IEEEtran}
\IEEEoverridecommandlockouts

\usepackage{cite}
\usepackage{amsmath,amssymb,amsfonts}
\usepackage{algorithmic}
\usepackage{graphicx}
\usepackage{textcomp}
\usepackage{xcolor}
\usepackage{booktabs}
\usepackage{array}

\begin{document}

\title{Improved Priority Exchange Server: Technical Briefing\\
{\footnotesize Milestone 1 --- Special Emphasis: Real-Time Systems, HSHL}}

\author{\IEEEauthorblockN{Sumon Boiddo}
\IEEEauthorblockA{\textit{BEng in Electronic Engineering} \\
\textit{Hochschule Hamm-Lippstadt}\\
Matriculation Number: 2220054 \\
Group B, Team 6}}

\maketitle

\begin{abstract}
This document is the Milestone~1 technical briefing for the seminar topic
\emph{Improved Priority Exchange Server} (IPE). It explains the problem that
motivates IPE, its core idea of reusing precomputed idle times, the formal
definition with its replenishment mechanism, the scheduling rules inherited
from the Dynamic Priority Exchange server, the schedulability condition, the
underlying assumptions, the run-time and memory overhead, the relevance of
IPE for embedded systems, and the open questions for further study. The
primary reference is Buttazzo, \textit{Hard Real-Time
Computing Systems}~\cite{buttazzo}, Chapter~6.
\end{abstract}

% ============================================================
\section{Problem Statement and Motivation}
% ============================================================

\begin{itemize}
  \item A real-time system runs two kinds of tasks together. \emph{Periodic}
        tasks arrive at fixed, regular intervals that are known in advance ---
        for example a temperature sensor reading data every second, or a task
        released every 10~ms --- and they have \emph{hard} deadlines that must
        never be missed~\cite{buttazzo}.
  \item \emph{Aperiodic} tasks arrive at random, unpredictable times --- for
        example an emergency such as an alarm or a sudden external request.
        They have soft deadlines and should be served as fast as possible.
  \item There are two simple ways to handle them, but both are bad. If the
        aperiodic task runs only in the \emph{background}, it waits until all
        periodic tasks finish, so when the system is busy its response time
        becomes very poor. If instead it runs at the \emph{highest priority},
        it keeps delaying the periodic tasks, which can make them miss a hard
        deadline --- and in a hard real-time system that means \textbf{system
        failure}.
  \item A better foundation is the \emph{Earliest Deadline First} (EDF)
        scheduler. With many tasks, Rate Monotonic (RM) can safely use only
        about $69\%$ of the CPU, but EDF can use up to $100\%$. Because more
        CPU is available under EDF, a larger aperiodic server can be used, so
        aperiodic requests receive better service~\cite{buttazzo}.
  \item Two EDF servers already exist. The \emph{Dynamic Priority Exchange}
        (DPE) server does not waste its free time: when no aperiodic request
        is waiting, it lends its time to a periodic task and saves an equal
        amount to reuse later~\cite{spuri1996}. The \emph{Earliest Deadline
        Late} (EDL) server is \emph{optimal} --- it gives the fastest possible
        aperiodic response --- but it must recompute the schedule at run time,
        which makes it too expensive for practical use~\cite{spuri1994}.
  \item \textbf{Core problem:} how can a server reach the near-optimal
        responsiveness of EDL \emph{without} paying EDL's heavy run-time cost?
        The Improved Priority Exchange (IPE) server answers this by computing
        the EDL idle times \emph{off-line, before running}, and reusing them
        inside a DPE-style server --- keeping EDL's good result while removing
        the expensive part.
\end{itemize}

% ============================================================
\section{Core Technical Contribution}
% ============================================================

\begin{itemize}
  \item IPE is a dynamic-priority (EDF) aperiodic server. Its main
        contribution is the way it \emph{combines two existing servers}: it
        uses the working method of the DPE server (the priority-exchange
        mechanism) together with the idle times of the optimal EDL
        server~\cite{spuri1996}.
  \item The key idea is to calculate the EDL idle times \emph{off-line, before
        the system runs}, and store them. This is comparable to preparing a
        complete schedule in advance: once the schedule is ready, no further
        decisions or calculations are needed while the work is carried out. In
        the same way, IPE performs its timing calculation beforehand, so that
        at run time it can serve an aperiodic (emergency) task very quickly,
        without expensive recomputation.
  \item Because the heavy work is done off-line, IPE avoids the high run-time
        cost of the EDL server while still serving aperiodic requests fast.
  \item IPE keeps the priority-exchange mechanism of DPE: when no aperiodic
        request is waiting, the server lends its free time to a periodic task
        and saves an equal amount of capacity to reuse later, so no CPU time
        is wasted.
  \item As a result, IPE behaves almost the same as the optimal EDL server ---
        it is \emph{nearly optimal} --- but at a much lower, practical
        run-time cost~\cite{buttazzo}.
\end{itemize}

% ============================================================
\section{Formal Definition}
% ============================================================

IPE is built on the idea of \emph{idle time} --- the moments when the CPU is
free. This is described by the availability function $\omega(t)$:
\begin{equation}
  \omega(t) =
  \begin{cases}
    1 & \text{if the CPU is idle (free) at time } t \\
    0 & \text{if the CPU is busy at time } t .
  \end{cases}
\end{equation}
The total idle time in an interval $[t_1,t_2]$ is the sum of all the free
moments in that interval:
\begin{equation}
  \Omega(t_1,t_2) = \int_{t_1}^{t_2} \omega(t)\, dt .
\end{equation}

IPE uses the idle times of the EDL schedule, which has a key property: among
all scheduling methods, EDL leaves the \emph{maximum} free time available as
early as possible~\cite{cc1989}. Because the free time comes early, an
arriving aperiodic request is served sooner, its waiting time is reduced, and
the periodic deadlines stay safe.

The pattern of EDL free times repeats every hyperperiod
$H = \mathrm{lcm}(T_1,\dots,T_n)$ (the least common multiple of all task
periods). The free times are therefore calculated only once, off-line, for a
single hyperperiod, and stored in two arrays:
\begin{itemize}
  \item $E = (e_0, e_1, \dots, e_p)$ --- the times at which the free periods
        occur (``when'');
  \item $D = (\Delta_0, \Delta_1, \dots, \Delta_p)$ --- the lengths of those
        free periods (``how long'').
\end{itemize}

The IPE server is then defined by the following rules~\cite{spuri1996}:
\begin{itemize}
  \item the server has an aperiodic capacity, initially set to $0$;
  \item at each instant $t = e_i + kH$ (with $k = 0, 1, 2, \dots$), the
        capacity is replenished by $\Delta_i$ units --- exactly the stored
        free time;
  \item the server always runs at the highest priority in the system;
  \item all other rules (serving aperiodic requests, executing periodic tasks,
        and the exchange and saving of capacity) are the same as in the DPE
        server.
\end{itemize}

\noindent\textbf{Schedulability condition.} The periodic tasks together with
the IPE server are schedulable (no hard deadline is ever missed) if and only
if
\begin{equation}
  U_p \leq 1 ,
\end{equation}
which is necessary because the CPU cannot do more than $100\%$ of work. The
server automatically allocates the remaining bandwidth $1 - U_p$ to the
aperiodic requests~\cite{spuri1996}.

% ============================================================
\section{Algorithm and Exchange/Replenishment Rules}
% ============================================================

This section describes, step by step, how the IPE server works.

\begin{itemize}
  \item \textbf{Replenishment.} The server starts with capacity $0$. At the
        precomputed times $t = e_i + kH$, it is refilled (replenished) by
        $\Delta_i$ units --- exactly the EDL free time stored in the arrays
        $E$ and $D$.
  \item \textbf{Exchange (no aperiodic request waiting).} When the server has
        capacity but no aperiodic task is waiting, the time is not wasted: a
        periodic task runs now, and an equal amount of capacity is saved
        (exchanged), so the server can use it later when an aperiodic request
        arrives. In this way no CPU time is wasted.
  \item \textbf{Consumption (aperiodic request arrives).} When an aperiodic
        request arrives and the server has capacity, it is served immediately,
        because the server always runs at the highest priority. As the request
        runs, the capacity decreases --- each unit of execution uses one unit
        of capacity. When the capacity reaches $0$, the server is exhausted
        and can no longer serve aperiodic requests; the periodic tasks run
        again, and the server waits for the next replenishment.
  \item \textbf{Spare-time reclaiming.} When a periodic task finishes earlier
        than its planned (worst-case) time --- for example, it was given
        10~units but finished in 5 --- the leftover time is not wasted. It is
        added to the server's capacity so an aperiodic request can use it.
        This is safe and breaks no periodic deadline, because that time was
        already reserved for the periodic task~\cite{buttazzo}.
\end{itemize}

% ============================================================
\section{Assumptions and Scope Limitations}
% ============================================================

\begin{itemize}
  \item The complete periodic task set must be known in advance (off-line),
        before the system runs --- each periodic task with its period and its
        computation time. This is required so that IPE can precompute the
        free-time arrays $E$ and $D$ off-line.
  \item The system uses a single processor (one CPU). The basic IPE model is
        defined for one CPU, not for multiprocessors.
  \item All tasks are fully preemptive --- a running task can be paused and
        continued later. This matters because an aperiodic request must be
        able to interrupt a running periodic task and run immediately.
  \item For the aperiodic tasks, the arrival time is unknown (they come
        suddenly and randomly), but their computation time (worst-case
        execution time, WCET) is known.
  \item Each periodic task has a period $T_i$, a computation time $C_i$, and a
        relative deadline equal to its period ($D_i = T_i$). All periodic
        tasks are assumed to be activated together at time $t = 0$.
  \item Some of these assumptions can be relaxed with a modified analysis ---
        for example arbitrary task phasing, deadlines different from periods,
        and shared resources (handled by the Stack Resource Policy). However,
        the single-processor and off-line-known task-set assumptions are
        fundamental to the basic IPE method~\cite{buttazzo}.
\end{itemize}

% ============================================================
\section{Runtime and Overhead}
% ============================================================

\begin{itemize}
  \item The heavy calculation --- the $E$ and $D$ arrays --- is done off-line,
        before the system runs. So at run time IPE only uses the stored result
        and does not recalculate anything. This makes its run-time cost low and
        fast.
  \item Compared to the EDL server, which must recompute the schedule at run
        time (high cost, slow), IPE moves all that heavy work off-line. As a
        result, its run-time overhead is low --- comparable to the DPE
        server~\cite{buttazzo}.
  \item The main cost of IPE is \emph{memory}. IPE must store the $E$ and $D$
        arrays, and their size depends on the hyperperiod
        $H = \mathrm{lcm}(T_1,\dots,T_n)$.
  \item If the periods fit nicely together (harmonic), $H$ is small, so the
        tables are small and need little memory. But if the periods do not fit
        nicely (non-harmonic), $H$ becomes very large, the stored tables become
        very large, and a lot of memory is needed. This is the main practical
        limitation of IPE~\cite{buttazzo}.
\end{itemize}

% ============================================================
\section{Relevance for Embedded Systems}
% ============================================================

\begin{itemize}
  \item IPE is a good choice for embedded systems where the periodic task set
        is fixed and known in advance (static). Because everything is known
        beforehand, IPE can precompute the $E$ and $D$ arrays off-line and then
        serve aperiodic requests very fast at run time.
  \item It works best when the periods are nicely matched (harmonic), so the
        hyperperiod stays small and the stored tables need little memory; when
        the platform supports an EDF scheduler; and when the system needs fast
        response to aperiodic events while still guaranteeing all hard periodic
        deadlines.
  \item IPE is a poor choice when the tasks are not known in advance --- that
        is, in a dynamic system where the task set changes while running ---
        because then the $E$ and $D$ arrays cannot be precomputed.
  \item It is also a poor choice when the periods do not fit nicely
        (non-harmonic), because the hyperperiod becomes huge and the stored
        tables need too much memory; and when the platform supports only fixed
        priority (RM) with no EDF scheduler.
\end{itemize}

% ============================================================
\section{Open Questions}
% ============================================================

\begin{itemize}
  \item \textbf{Memory limitation.} How large can the memory cost of the $E$
        and $D$ arrays actually become for realistic non-harmonic task sets,
        and can it be reduced or bounded?
  \item \textbf{Comparison with EDL.} How close is IPE to the optimal EDL
        server in practice, and in which load conditions does EDL still perform
        slightly better?
  \item \textbf{Comparison with other servers.} How does IPE compare with
        simpler servers (such as the Total Bandwidth Server) in terms of
        responsiveness versus implementation simplicity?
  \item \textbf{Relaxing the assumptions.} What happens if the base assumptions
        are relaxed --- for example a changing (dynamic) task set, deadlines
        different from periods, or shared resources?
\end{itemize}

% ============================================================
% References
% ============================================================

\begin{thebibliography}{9}

\bibitem{buttazzo}
G.~C.~Buttazzo,
\textit{Hard Real-Time Computing Systems: Predictable Scheduling Algorithms
and Applications}, 3rd~ed.
Springer, 2011, ch.~6 (Dynamic Priority Servers).

\bibitem{spuri1994}
M.~Spuri and G.~Buttazzo,
``Efficient Aperiodic Service under Earliest Deadline Scheduling,''
in \textit{Proc. IEEE Real-Time Systems Symposium (RTSS)}, 1994.

\bibitem{spuri1996}
M.~Spuri and G.~Buttazzo,
``Scheduling Aperiodic Tasks in Dynamic Priority Systems,''
\textit{Real-Time Systems}, vol.~10, no.~2, March 1996.

\bibitem{lss1987}
J.~Lehoczky, L.~Sha, and J.~Strosnider,
``Enhanced Aperiodic Responsiveness in Hard Real-Time Environments,''
in \textit{Proc. IEEE Real-Time Systems Symposium (RTSS)}, 1987.

\bibitem{cc1989}
H.~Chetto and M.~Chetto,
``Some Results of the Earliest Deadline Scheduling Algorithm,''
\textit{IEEE Trans. Software Engineering}, vol.~15, no.~10, 1989.

\bibitem{ai_usage_m1}
S.~Boiddo,
\textit{AI Usage Protocol, Milestone 1: Improved Priority Exchange Server},
HSHL, 2026. Protocol version 0.2; JSON and BibTeX files submitted with this
milestone.

\end{thebibliography}

\end{document}
